\subsection{阿贝尔扩张与阿贝尔闭包}

定义 1. --- 若域 K 的扩张 E 是伽罗瓦扩张且其伽罗瓦群是交换的，则称 E 为阿贝尔扩张。

由于交换群的每个子群都是正规子群，V 第68页的推论4表明阿贝尔扩张的每个子扩张都是阿贝尔的。

命题 1. --- 设 E 是 K 的伽罗瓦扩张，Γ 是其伽罗瓦群。设 Δ 为 Γ 的导出群（I，第 10 页，定义 4），F 为 Δ 的不变域。对于 E 的子扩张 L，其为阿贝尔扩张当且仅当 L 包含于 F 中。

首先我们注意到，F 也是闭包 \(\bar{\Delta}\) （ \(\Delta\) 在 \(\Gamma\) 中的）的不变量域，并且 \(\bar{\Delta}\) 是 \(\Gamma\) 的闭正规子群。根据 V，第68页，推论4，F 因此是 K 的伽罗瓦扩张。进一步地，F 在 K 上的伽罗瓦群同构于 \(\Gamma/\bar{\Delta}\) ，从而是交换的。因此 F 的每个子扩张都是阿贝尔的。反之，设 L 是包含于 E 中的 K 的一个阿贝尔扩张，并令 \(\Pi\) 为 E 在 L 上的伽罗瓦群。由于 L 是伽罗瓦的， \(\Pi\) 是 \(\Gamma\) 的一个正规子群，并且 L 在 K 上的伽罗瓦群同构于 \(\Gamma/\Pi\) (V，第68页，推论4)。因此 \(\Gamma/\Pi\) 是交换的，并且 \(\Pi\) 包含 \(\Delta\) ，从而 \(L \subset F\) 。

推论. --- 设 E 是 K 的一个扩域，并设 \((\mathrm{E}_{i})_{i\in I}\) 是 E 的一族子扩域，使得 \(\mathrm{E}=\mathrm{K}\left(\bigcup_{i\in I}\mathrm{E}_{i}\right)\) 。假设每个扩域

\(E_{i}\) 都是阿贝尔的，则 E 亦然。

首先，E 是 K 的伽罗瓦扩张（V，第57页，命题1）。如果域 F 按命题1中的方式定义，那么我们有 \(E_{i} \subset F\) （对所有 \(i \in I\) ），因此 E = F。

域 K 的扩张 E 被称为 K 的阿贝尔闭包，如果它是 K 的阿贝尔扩张，并且 K 的每个阿贝尔扩张都同构于 E 的一个子扩张。命题1蕴含 K 的阿贝尔闭包的存在性：因为设 \(K_{s}\) 为 K 的一个可分闭包，其伽罗瓦群为 \(\Gamma\) ，并且令 \((\overline{\Gamma},\overline{\Gamma})\) 为 \(\Gamma\) 的导出群的闭包；记 \(K_{ab}\) 为 \((\overline{\Gamma},\overline{\Gamma})\) 的不变量域；由于 K 的每个可分代数扩张都同构于 \(K_{s}\) 的某个子扩张（V，第45页，推论），命题1表明 \(K_{ab}\) 是 K 的阿贝尔闭包。 \(K_{ab}\) 在 K 上的伽罗瓦群典范同构于 \(\Gamma/(\overline{\Gamma},\overline{\Gamma})\) 。现在让我们证明阿贝尔闭包的唯一性：设 E， \(E'\) 是 K 的两个阿贝尔闭包；根据定义，存在 K-同态 \(u:E\to E'\) 和 \(v:E'\to E\) ，且命题1（V，第52页）蕴含 \(v(u(E)) = E\) 和 \(u(v(E')) = E'\) ，因此 \(u\) 是 \(E\) 到 \(E'\) 上的一个 K-同构。任何其他从 \(E\) 到 \(E'\) 上的 K-同构都具有形式 \(u_1 = \sigma_0 \circ u\) ，其中 \(\sigma_0 \in \operatorname{Gal}(E'/K)\) ；由于 \(\operatorname{Gal}(E'/K)\) 是交换的，同构 \(\sigma \mapsto u \circ \sigma \circ u^{-1}\) （从 \(\operatorname{Gal}(E/K)\) 到 \(\operatorname{Gal}(E'/K)\) 上）与 \(u\) 无关；它被称为 \(\operatorname{Gal}(E/K)\) 到 \(\operatorname{Gal}(E'/K)\) 上的典范同构。

